% !TeX root = part5.tex
\documentclass[../final.tex]{subfiles}
\begin{document}
\ifSubfilesClassLoaded{\setcounter{chapter}{2}\setcounter{section}{6}}{}

\chapter{Вырожденный бозе-газ}
\rsubtitle{Конденсация Бозе--Эйнштейна и теплоёмкость}

\Needspace{7\baselineskip}
\section{Насыщение возбуждённых состояний}
\subsection{Распределение Бозе--Эйнштейна}
Рассматриваем \rassumption{трёхмерный идеальный газ с сохраняющимся числом
частиц}, спектром $\varepsilon=p^2/(2m)$ и нулевой энергией основного состояния.
Температура измеряется в энергетических единицах: $T=k_BT_{\mathrm K}$.
Среднее заполнение одного состояния равно
\begin{equation*}
 \overline n_i=\frac1{e^{(\varepsilon_i-\mu)/T}-1}.
\end{equation*}
Положительность заполнения требует $\mu<0$ в конечной системе;
в термодинамическом пределе допускается $\mu=0$.
Плотность возбуждённых состояний и их суммарное заполнение имеют вид
\begin{align*}
 D(\varepsilon)&=AV\sqrt\varepsilon,\qquad
 A=\frac{g(2m)^{3/2}}{4\pi^2\hbar^3},\\
 N_{\mathrm{ex}}&=AV\int_0^\infty
 \frac{\sqrt\varepsilon\,d\varepsilon}{e^{(\varepsilon-\mu)/T}-1}
 =AVT^{3/2}I(x),\qquad x=\frac\mu T\leq0,\\
 I(x)&=\int_0^\infty\frac{\sqrt z\,dz}{e^{z-x}-1}.
\end{align*}
Здесь $g$ --- кратность внутреннего вырождения, $z=\varepsilon/T$.
До конденсации основное состояние заполнено немакроскопически и
$N_{\mathrm{ex}}\simeq N$.

\subsection{Максимальное число возбуждённых частиц}
При $x<0$ функция $I$ монотонно возрастает:
\begin{equation*}
 I'(x)=\int_0^\infty
 \frac{\sqrt z\,e^{z-x}}{(e^{z-x}-1)^2}\,dz>0.
\end{equation*}
Максимум достигается при $\mu\to0^{-}$, когда
$I(0)=\Gamma(3/2)\zeta(3/2)$:
\begin{equation*}
 N_{\mathrm{ex}}^{\max}(T,V)=
 AVT^{3/2}\Gamma\!\left(\frac32\right)\zeta\!\left(\frac32\right).
\end{equation*}
При $N=\mathrm{const}$ уменьшение $T$ или $V$ увеличивает $N/(VT^{3/2})$.
Химический потенциал приближается к нулю, после чего дальнейшее увеличение
этого отношения уже невозможно компенсировать изменением $\mu$.

\clearpage
\section{Конденсация Бозе--Эйнштейна}
\subsection{Критическая температура и критический объём}
\rconcept{Температура конденсации Бозе--Эйнштейна} $T_c=T_{\text{БЭ}}$
определяется равенством $N=N_{\mathrm{ex}}^{\max}(T_c,V)$:
\begin{align*}
 T_c&=\left[
 \frac{N}{AV\Gamma(3/2)\zeta(3/2)}\right]^{2/3}
 =\frac{2\pi\hbar^2}{m}
 \left[\frac{N}{gV\zeta(3/2)}\right]^{2/3},\\
 V_c(T)&=\frac{N}{AT^{3/2}\Gamma(3/2)\zeta(3/2)}.
\end{align*}
Использованы $\Gamma(3/2)=\sqrt\pi/2$ и $\zeta(3/2)\simeq2.612$.
Для температуры в кельвинах правую часть формулы для $T_c$ делят на $k_B$.
При фиксированном $V$ конденсация происходит при $T<T_c$;
при фиксированном $T$ --- при сжатии до $V<V_c$.

\subsection{Основное состояние и интегрируемая особенность}
При $\mu=0$ заполнение состояния с $\varepsilon=0$ формально расходится.
Однако интеграл по возбуждённым состояниям остаётся конечным:
\begin{equation*}
 \int_0^\delta\frac{\sqrt\varepsilon\,d\varepsilon}{e^{\varepsilon/T}-1}
 \simeq T\int_0^\delta\frac{d\varepsilon}{\sqrt\varepsilon}
 =2T\sqrt\delta\ \longrightarrow\ 0\quad(\delta\to0).
\end{equation*}
Макроскопическое заполнение отдельного основного состояния учитывается
отдельно от непрерывной плотности состояний:
\begin{align*}
 N&=N_0+N_{\mathrm{ex}},\qquad
 N_{\mathrm{ex}}=AVT^{3/2}\Gamma(3/2)\zeta(3/2),\qquad \mu=0,\\
 \frac{N_0}{N}&=1-\left(\frac{T}{T_c}\right)^{3/2}\quad(T\leq T_c),\qquad
 \frac{N_0}{N}=1-\frac{V}{V_c}\quad(V\leq V_c).
\end{align*}
При $T>T_c$ или $V>V_c$ доля конденсата равна нулю.
В конечной системе большое $N_0$ соответствует малому отрицательному
$\mu$: при невырожденном основном уровне
$N_0=(e^{-\mu/T}-1)^{-1}\simeq T/(-\mu)$.

\begin{rbox}[unbreakable,left=18pt,right=18pt,top=14pt,bottom=10pt]{[]}
\centering
\captionsetup{font=footnotesize,hypcap=false,skip=7pt}
\begin{tikzpicture}[>=Stealth,font=\small,line width=0.8pt]
 \foreach \dx/\xl in {0/{$T/T_c$},6.2/{$V/V_c$}} {
  \begin{scope}[xshift=\dx cm]
   \draw[->] (0,0)--(4.8,0) node[right] {\xl};
   \draw[->] (0,0)--(0,2.7) node[above] {$N_0/N$};
   \node[left] at (0,2.2) {$1$};
   \draw[dotted] (3,0)--(3,2.3);
   \node[below] at (3,0) {$1$};
   \draw[p3,thick] (3,0)--(4.4,0);
  \end{scope}
 }
 \draw[p3,thick,domain=0:3,samples=65]
  plot (\x,{2.2*(1-(\x/3)^1.5)});
 \draw[p3,thick] (6.2,2.2)--(9.2,0);
\end{tikzpicture}
\captionof{figure}{Доля частиц в конденсате: слева при охлаждении с постоянными
 $N,V$, справа при сжатии с постоянными $N,T$. На границе перехода $N_0=0$;
 ниже критической температуры или критического объёма $N_0$ становится макроскопическим.}
\end{rbox}

\clearpage
\section{Бозе-интегралы и специальные функции}
\subsection{Интеграл при нулевом химическом потенциале}
Для $\alpha>1$ разложим знаменатель в геометрический ряд:
\begin{align*}
 \frac1{e^z-1}&=\sum_{n=1}^{\infty}e^{-nz},\\
 \int_0^\infty\frac{z^{\alpha-1}\,dz}{e^z-1}
 &=\sum_{n=1}^{\infty}\int_0^\infty z^{\alpha-1}e^{-nz}\,dz
 =\Gamma(\alpha)\sum_{n=1}^{\infty}\frac1{n^\alpha}
 =\Gamma(\alpha)\zeta(\alpha).
\end{align*}
Здесь $\Gamma(\alpha)=\int_0^\infty u^{\alpha-1}e^{-u}du$,
а $\zeta(\alpha)=\sum_{n=1}^\infty n^{-\alpha}$ при $\alpha>1$.
Аналогично для ферми-газа
\begin{equation*}
 \int_0^\infty\frac{z^{\alpha-1}\,dz}{e^z+1}
 =\Gamma(\alpha)\sum_{n=1}^{\infty}\frac{(-1)^{n-1}}{n^\alpha}
 =\Gamma(\alpha)(1-2^{1-\alpha})\zeta(\alpha),\qquad \alpha>1.
\end{equation*}
Условие $\alpha>1$ для бозе-интеграла связано с поведением
$z^{\alpha-2}$ при $z\to0$.

\subsection{Ненулевой химический потенциал}
Введём функцию, используемую в конспекте:
\begin{align*}
 F_\alpha(x)&=\sum_{n=1}^{\infty}\frac{e^{nx}}{n^\alpha},\qquad x<0,\\
 \int_0^\infty\frac{z^{\alpha-1}\,dz}{e^{z-x}-1}
 &=\Gamma(\alpha)F_\alpha(x),\qquad \alpha>0.
\end{align*}
Она совпадает с полилогарифмом $\operatorname{Li}_\alpha(e^x)$.
При $x<0$ ряд допускает почленное дифференцирование:
\begin{equation*}
 F_\alpha'(x)=F_{\alpha-1}(x),\qquad
 F_\alpha(0)=\zeta(\alpha)\quad(\alpha>1).
\end{equation*}
Следовательно,
\begin{equation*}
 I(x)=\Gamma(3/2)F_{3/2}(x),\qquad
 I'(x)=\Gamma(3/2)F_{1/2}(x).
\end{equation*}
Сам $I(0)$ конечен, но $I'(x)$ расходится при $x\to0^{-}$.
Разложение в обычный ряд Тейлора около нуля допустимо для первого
поправочного члена только при $\alpha>2$:
\begin{equation*}
 F_\alpha(x)=\zeta(\alpha)+\zeta(\alpha-1)x+o(|x|),\qquad\alpha>2.
\end{equation*}
При $1<\alpha<2$ главная поправка содержит дробную степень $(-x)^{\alpha-1}$.

\clearpage
\section{Неаналитическое поведение вблизи перехода}
\subsection{Выделение вклада малых энергий}
Положим $y=-\mu/T>0$ и рассмотрим $0<a<1$:
\begin{equation*}
 b_a(y)=\int_0^\infty\frac{u^a\,du}{e^{u+y}-1}
 =\Gamma(a+1)F_{a+1}(-y).
\end{equation*}
Функция $b_a(0)$ конечна, но её производная при $y\to0^+$ определяется
областью малых $u$ и расходится. В этой области $e^{u+y}-1\simeq u+y$,
поэтому главный сингулярный член равен
\begin{align*}
 b_a'(y)&=-\int_0^\infty
 \frac{u^a e^{u+y}}{(e^{u+y}-1)^2}\,du
 \sim-\int_0^\infty\frac{u^a\,du}{(u+y)^2}\\
 &=-y^{a-1}\int_0^\infty\frac{v^a\,dv}{(1+v)^2}
 =-y^{a-1}\Gamma(a+1)\Gamma(1-a).
\end{align*}
Подстановка $u=yv$ выделила всю зависимость от малого параметра.
Последний интеграл вычисляется через бета-функцию:
\begin{equation*}
 B(p,q)=\int_0^\infty\frac{v^{p-1}\,dv}{(1+v)^{p+q}}
 =\frac{\Gamma(p)\Gamma(q)}{\Gamma(p+q)},\qquad p,q>0.
\end{equation*}
Интегрируя главный член производной и используя
$\Gamma(1-a)=-a\Gamma(-a)$, получаем
\begin{align*}
 b_a(y)&=b_a(0)-\frac{\Gamma(a+1)\Gamma(1-a)}{a}\,y^a+o(y^a)\\
 &=\Gamma(a+1)\bigl[\zeta(a+1)+\Gamma(-a)y^a+o(y^a)\bigr],\\
 F_\alpha(x)&=\zeta(\alpha)+\Gamma(1-\alpha)(-x)^{\alpha-1}
 +o\!\left((-x)^{\alpha-1}\right),\qquad1<\alpha<2.
\end{align*}
При замене $y=z^{1/a}$ главный член $b_a(z^{1/a})-b_a(0)$ становится
линейным по $z$. Для $a=1/2$ это подстановка $y=z^2$, использованная в лекции.

\subsection{Случай трёхмерного газа}
Для $\alpha=3/2$ и $\Gamma(-1/2)=-2\sqrt\pi$:
\begin{align*}
 F_{3/2}(-y)&=\zeta(3/2)-2\sqrt\pi\sqrt y+O(y),\\
 F_{1/2}(-y)&=\frac{\sqrt\pi}{\sqrt y}+O(1),\qquad y\to0^+.
\end{align*}
Именно корневая особенность определяет поведение $\mu$ и теплоёмкости
вблизи конденсации.

\clearpage
\subsection{Химический потенциал при фиксированных числе частиц и объёме}
В нормальной фазе $T>T_c$ уравнение для числа частиц имеет вид
\begin{equation*}
 N=AVT^{3/2}\Gamma(3/2)F_{3/2}(x),\qquad x=\frac\mu T<0.
\end{equation*}
Используем это же равенство при $T=T_c$, где $x=0$:
\begin{equation*}
 F_{3/2}(x)=\zeta(3/2)\left(\frac{T_c}{T}\right)^{3/2}.
\end{equation*}
Для $\tau=(T-T_c)/T_c\to0^+$ из корневой асимптотики следует
\begin{align*}
 2\sqrt\pi\sqrt{-x}
 &=\zeta(3/2)\left[1-(1+\tau)^{-3/2}\right]+O(\tau^2)\\
 &=\frac32\zeta(3/2)\tau+O(\tau^2),\\
 \frac\mu T=x&=-\frac{9\zeta^2(3/2)}{16\pi}\tau^2+O(\tau^3).
\end{align*}
Таким образом, $\mu$ приближается к нулю квадратично.
Тот же вывод получается при дифференцировании уравнения для $N$:
\begin{align*}
 0&=\frac32F_{3/2}(x)+TF_{1/2}(x)\frac{dx}{dT},\\
 \frac{dx}{dT}&=-\frac{3}{2T}\frac{F_{3/2}(x)}{F_{1/2}(x)},\qquad
 \left.\frac{dx}{dT}\right|_{T\to T_c^+}=0,\\
 \left.\frac{d^2x}{dT^2}\right|_{T\to T_c^+}
 &=-\frac{9\zeta^2(3/2)}{8\pi T_c^2}.
\end{align*}
В фазе с конденсатом $\mu=0$ при $T\leq T_c$ в термодинамическом пределе.
Первая производная $\mu$ непрерывна и равна нулю на границе,
а вторая производная имеет конечный скачок.

\begin{rbox}[unbreakable,left=18pt,right=18pt,top=14pt,bottom=10pt]{[]}
\centering
\captionsetup{font=footnotesize,hypcap=false,skip=7pt}
\begin{tikzpicture}[>=Stealth,font=\small,line width=0.8pt,x=1.3cm,y=2cm]
 \foreach \dx/\xl in {0/{$T/T_c$},6.2/{$V/V_c$}} {
  \begin{scope}[xshift=\dx cm]
   \draw[->] (0,0)--(3.65,0) node[right] {\xl};
   \draw[->] (0,-1.15)--(0,0.4) node[above] {$\mu/T$};
   \draw[dotted] (1,0)--(1,-1.05);
   \node[above] at (1,0) {$1$};
   \draw[p3,very thick] (0.08,0)--(1,0);
   \node[p3,above] at (0.52,0.1) {$\mu=0$};
   \node[p3,below] at (2.75,0) {$\mu<0$};
  \end{scope}
 }
 \draw[p3,thick] plot file {\subfix{../image/bose_mu.dat}};
 \begin{scope}[xshift=6.2cm]
  \draw[p3,thick] plot file {\subfix{../image/bose_mu_volume.dat}};
 \end{scope}
\end{tikzpicture}
\captionof{figure}{Химический потенциал: слева при постоянных $N,V$, справа
 при постоянных $N,T$. Кривые вычислены из уравнения для числа частиц.
 При $T\leq T_c$ или $V\leq V_c$ имеется конденсат и $\mu=0$;
 со стороны нормальной фазы графики подходят к нулю с горизонтальной касательной.}
\end{rbox}

\clearpage
\section{Энергия и теплоёмкость}
\subsection{Внутренняя энергия}
Основное состояние с $\varepsilon_0=0$ не вносит вклада в энергию.
Энергия определяется возбуждёнными частицами:
\begin{align*}
 E&=AV\int_0^\infty
 \frac{\varepsilon^{3/2}\,d\varepsilon}{e^{(\varepsilon-\mu)/T}-1}\\
 &=AVT^{5/2}\Gamma(5/2)
 \begin{cases}
  F_{5/2}(\mu/T),&T>T_c,\\
  \zeta(5/2),&T\leq T_c.
 \end{cases}
\end{align*}
Здесь $\Gamma(5/2)=3\sqrt\pi/4$, $\zeta(5/2)\simeq1.341$.
При $T<T_c$ число возбуждённых частиц пропорционально $T^{3/2}$,
а характерная энергия одной такой частицы --- $T$, поэтому $E\propto T^{5/2}$.

\subsection{Производная энергии при постоянных числе частиц и объёме}
Теплоёмкость в принятых энергетических единицах определяется как
$C_V=(\partial E/\partial T)_{N,V}$.
В конденсированной фазе $\mu=0$, поэтому
\begin{align*}
 C_V^-&=\frac52 AVT^{3/2}\Gamma(5/2)\zeta(5/2),\\
 \frac{C_V^-}{N}&=\frac{15}{4}\frac{\zeta(5/2)}{\zeta(3/2)}
 \left(\frac{T}{T_c}\right)^{3/2}.
\end{align*}
Знак $-$ обозначает ветвь $T<T_c$. В нормальной фазе
$x=\mu/T$ зависит от температуры:
\begin{align*}
 C_V^+&=AV\Gamma(5/2)
 \left[\frac52T^{3/2}F_{5/2}(x)
 +T^{5/2}F_{3/2}(x)\frac{dx}{dT}\right]\\
 &=AVT^{3/2}\Gamma(5/2)
 \left[\frac52F_{5/2}(x)-\frac32\frac{F_{3/2}^2(x)}{F_{1/2}(x)}\right],\\
 \frac{C_V^+}{N}&=\frac{15}{4}\frac{F_{5/2}(x)}{F_{3/2}(x)}
 -\frac94\frac{F_{3/2}(x)}{F_{1/2}(x)},\qquad T>T_c.
\end{align*}
При $T\gg T_c$ имеем $e^x\ll1$ и $F_\alpha(x)\simeq e^x$.
Тогда $E\simeq3NT/2$ и $C_V\to3N/2$: восстанавливается
классический результат для идеального газа.
Теплоёмкость по температуре в кельвинах равна $k_BC_V$;
безразмерная величина $C_V/N$ соответствует обычному $C_{V,\mathrm K}/(Nk_B)$.

\clearpage
\subsection{Непрерывность теплоёмкости на границе перехода}
При $T\to T_c^+$ имеем $x\to0^{-}$, $dx/dT\to0$,
$F_{5/2}(x)\to\zeta(5/2)$ и $F_{1/2}(x)\to\infty$.
Следовательно,
\begin{equation*}
 C_V(T_c^+)=C_V(T_c^-)
 =\frac{15}{4}N\frac{\zeta(5/2)}{\zeta(3/2)}\simeq1.926N.
\end{equation*}
Теплоёмкость непрерывна, достигает максимума при $T=T_c$,
но её наклон меняется скачком.

\begin{rbox}[unbreakable,left=18pt,right=18pt,top=14pt,bottom=10pt]{[]}
\centering
\captionsetup{font=footnotesize,hypcap=false,skip=7pt}
\begin{tikzpicture}[>=Stealth,font=\small,line width=0.8pt,x=2.65cm,y=1.6cm]
 \draw[->] (0,0)--(4.25,0) node[right] {$T/T_c$};
 \draw[->] (0,0)--(0,2.45) node[above] {$C_V/N$};
 \draw[dashed,red3] (0,1.5)--(4,1.5);
 \node[left] at (0,1.5) {$3/2$};
 \draw[dotted] (1,0)--(1,1.926);
 \node[below] at (1,0) {$1$};
 \draw[p3,thick,domain=0:1,samples=60] plot (\x,{1.9256717*\x^1.5});
 \draw[p3,thick] plot file {\subfix{../image/bose_cv.dat}};
 \fill[p3] (1,1.9256717) circle (1.7pt);
 \node[above,p3] at (1,1.98) {$1.926$};
 \node[p3,align=center] at (2.75,0.75) {нормальная\\фаза};
\end{tikzpicture}
\captionof{figure}{Теплоёмкость трёхмерного идеального бозе-газа при постоянных $N,V$.
 Ниже $T_c$ она растёт как $T^{3/2}$; выше $T_c$ уменьшается к классическому
 пределу $3N/2$. В точке перехода имеется излом без скачка самой теплоёмкости.}
\end{rbox}

\subsection{Скачок производной теплоёмкости}
При дифференцировании выражения для $C_V^+$ на границе перехода
выживает добавка, содержащая $d^2x/dT^2$:
\begin{align*}
 \left.\frac{dC_V^+}{dT}\right|_{T_c}
 -\left.\frac{dC_V^-}{dT}\right|_{T_c}
 &=AVT_c^{5/2}\Gamma(5/2)\zeta(3/2)
 \left.\frac{d^2x}{dT^2}\right|_{T_c^+}\\
 &=-\frac{27}{16\pi}\frac{N}{T_c}\zeta^2(3/2)
 \simeq-3.666\frac{N}{T_c}.
\end{align*}
Произведение $F_{1/2}(x)(dx/dT)^2$ при этом стремится к нулю:
оно имеет порядок $T-T_c$. Односторонние наклоны равны
\begin{equation*}
 \left.\frac{dC_V^-}{dT}\right|_{T_c}
 =\frac{3C_V(T_c)}{2T_c}\simeq2.889\frac{N}{T_c},\qquad
 \left.\frac{dC_V^+}{dT}\right|_{T_c}\simeq-0.777\frac{N}{T_c}.
\end{equation*}

\end{document}
